\documentclass[11pt,a4paper]{article}
\usepackage[margin=2.5cm]{geometry}
\usepackage{longtable}
\usepackage{booktabs}
\usepackage[T1]{fontenc}

\title{Practica 5B: Pruebas de software sobre \texttt{ConjuntoOrdenado}}
\author{Juan Calle \and Juan Francisco Gonzalez}
\date{\today}

\begin{document}
\maketitle

\section{Objetivo}
La practica consiste en disenar y ejecutar pruebas unitarias sobre la clase \texttt{ConjuntoOrdenado}, aplicando tecnicas de caja negra y caja blanca, detectando los defectos existentes y corrigiendolos posteriormente.

\section{Entorno y material entregado}
Se ha creado un proyecto Maven llamado \texttt{practica5b} con \texttt{groupId = es.unican.is2}. El proyecto depende de la libreria \texttt{AbstractDataTypes} version \texttt{1.0}, donde se encuentra la interfaz \texttt{IConjuntoOrdenado}. La entrega incluye:
\begin{itemize}
\item La implementacion de \texttt{ConjuntoOrdenado}.
\item La clase de pruebas \texttt{ConjuntoOrdenadoTest}.
\item El fichero \texttt{pom.xml} configurado para ejecutar pruebas con JUnit y medir cobertura con JaCoCo.
\item Este informe en \LaTeX.
\end{itemize}

\section{Diseno de pruebas de caja negra}
El diseno se ha realizado a partir de la especificacion de la interfaz \texttt{IConjuntoOrdenado}, usando particion equivalente y analisis de valores limite.

\subsection{Particiones equivalentes}
\begin{longtable}{p{2.8cm}p{5.2cm}p{6cm}}
\toprule
Metodo & Particiones & Casos implementados \\
\midrule
\endhead
\texttt{add} &
Elemento valido en conjunto vacio; elemento valido en conjunto no vacio; elemento duplicado; elemento nulo &
Insercion en vacio, insercion manteniendo orden, rechazo de duplicados, excepcion por nulo \\
\texttt{get} &
Indice valido; indice negativo; indice igual al tamano &
Acceso a posiciones validas, excepcion para $-1$, excepcion para \texttt{size()} \\
\texttt{remove} &
Indice valido; indice negativo; indice igual al tamano &
Eliminacion correcta de una posicion intermedia, excepcion para $-1$, excepcion para \texttt{size()} \\
\texttt{size} &
Conjunto vacio; conjunto no vacio &
Comprobacion de cardinalidad antes y despues de inserciones \\
\texttt{clear} &
Conjunto no vacio; conjunto vacio tras limpiar &
Vaciado del conjunto y comprobacion de inaccesibilidad posterior \\
\bottomrule
\end{longtable}

\subsection{Analisis de valores limite}
Los limites relevantes aparecen en los indices y en el cardinal del conjunto:
\begin{itemize}
\item Para \texttt{get} y \texttt{remove}: se prueban los limites invalidos \texttt{-1} y \texttt{size()}, y valores validos como \texttt{0} y una posicion intermedia.
\item Para \texttt{size}: se comprueban los valores limite \texttt{0}, \texttt{1}, \texttt{2} y \texttt{3} mediante distintas secuencias de insercion y borrado.
\item Para \texttt{add}: se prueban inserciones al principio, en medio y en un conjunto vacio, ademas del caso limite de duplicado.
\end{itemize}

\section{Pruebas de caja blanca}
Se han anadido pruebas orientadas a cubrir todos los caminos relevantes del metodo \texttt{add}, que era el unico metodo con logica de decision significativa:
\begin{itemize}
\item Camino con conjunto vacio.
\item Camino con conjunto no vacio y sin iteraciones del bucle.
\item Camino con conjunto no vacio y con iteraciones del bucle.
\item Camino de rechazo de duplicados.
\item Camino de excepcion por parametro nulo.
\end{itemize}
Ademas, se han cubierto tanto las ejecuciones normales como las excepcionales de \texttt{get} y \texttt{remove}.

\section{Ejecucion inicial y defectos detectados}
Las pruebas se ejecutaron primero sobre la version original defectuosa. La ejecucion inicial mostro fallos en cuatro pruebas:
\begin{itemize}
\item \texttt{testAddMantieneOrdenNatural}
\item \texttt{testAddNoInsertaDuplicados}
\item \texttt{testGetIndicesValidos}
\item \texttt{testRemoveIndiceValido}
\end{itemize}

El analisis de esos fallos permitio localizar dos defectos en el metodo \texttt{add}:
\begin{itemize}
\item La condicion del bucle estaba invertida y provocaba inserciones que rompian el orden natural.
\item No se comprobaba si el elemento ya existia, por lo que se admitian duplicados y siempre se devolvia \texttt{true}.
\end{itemize}

\section{Correccion realizada}
Se corrigio el metodo \texttt{add} para:
\begin{itemize}
\item Avanzar mientras el nuevo elemento sea mayor que el actual, conservando el orden ascendente.
\item Detectar si el elemento ya existe en la posicion de insercion y devolver \texttt{false} sin modificar el conjunto.
\end{itemize}

\section{Resultados finales}
Tras la correccion, la ejecucion de \texttt{mvn test} finaliza correctamente:
\begin{itemize}
\item 13 pruebas ejecutadas.
\item 0 fallos.
\item 0 errores.
\end{itemize}

La cobertura obtenida con JaCoCo es completa sobre la clase \texttt{ConjuntoOrdenado}:
\begin{itemize}
\item Instrucciones: 100\%.
\item Ramas: 100\%.
\item Lineas: 100\%.
\item Metodos: 100\%.
\end{itemize}

\section{Conclusion}
La bateria de pruebas diseniada permite detectar los defectos presentes en la implementacion original y verificar posteriormente la correccion de la clase. El proyecto entregado queda preparado para su ejecucion mediante Maven y cumple los requisitos funcionales especificados para la estructura de datos \texttt{ConjuntoOrdenado}.

\end{document}
